\chapter{Literature Survey}

\section{Early Automation} 

Early-stage automation relied on random, rule-based, and script-based strategies. These methods had notable limitations and do not provide high flexibility and scalability.

\begin{enumerate}
    \item \textbf{Random-based automation} Random testing is helpful for identifying bugs that structured automation may miss \cite{RandomTesting:inproceedings}. Applying a random operation on the GUI is the idea for the random-based automation. It is not the best tool due to its highly redundant and useless action performance.
    \item \textbf{Ruled-based automation} Ruled-based automation is planned automation with idea to apply rules on known UI structure \cite{RuledBasedAutomation}. It is a widely used testing strategy for exploring testing and bug detection
    \item \textbf{Script-based automation} To control GUI interactions, script-based automation uses code scripts.  Java binaries and the JVM are used by programs like jRapture \cite{ScriptBasedAutomation} to record and replay Java-based GUI sequences, allowing consistent execution by accurately replicating input sequences.
\end{enumerate}



\section{Shift Toward Intelligent GUI Automation} 
Machine Learning and Deep Learning based approaches motivate intelligent GUI automation. Furthermore, LLM and MLLM enhance this field with large prompt and visual input capabilities.\\ 
\textbf{GPT‑4V(ision) is a Generalist Web Agent, if Grounded (SeeAct):}\\
SeeAct \cite{SeeAct} is foundational work showing that the large multimodal models like GPT-4V can act as generalist web agents. It uses both visual and textual values to generate actions and use the Plawright environment for action executions \cite{SeeAct}. The paper achieve 51.1\% accuracy on online evaluation, which is 20-30\% improvement over only text-based methods \cite{SeeAct}. The paper also shows that grounding remains the main bottleneck and possible improvement suggestion on grounding \cite{SeeAct}. \\
\textbf{AdaptAgent: Adapting Multimodal Web Agents with Few‑Shot Learning from Human Demonstrations:}\\
This paper addresses a key limitation of SeeAct for unseen websites or domains. AdaptAgent proposes a few-shot learning based method over a new website or domain \cite{FewShotLearning:verma2024adaptagentadaptingmultimodalweb}. Experiments on standard benchmarks show that even with few shots, agents improve success rate by 3.36\% to 7.21\% \cite{FewShotLearning:verma2024adaptagentadaptingmultimodalweb}.
This suggests that few-shot adaptation is a promising, efficient way to make MLLM-based agents more robust across diverse websites, instead of relying solely on large-scale fine-tuning.\\
\textbf{Browsing Like Human: A Multimodal Web Agent with Experiential Fast-and-Slow Thinking (WebExperT):}\\
The goal of this latest effort is to improve web agents' human-like reasoning.  The authors present a dual-process planning paradigm: "fast" initial passes and "slow" deliberation when necessary \cite{Browsinglikehuman:luo2025browsing}. It also incorporates experience learning, whereby the agent considers its past mistakes and modifies its behaviour going forward \cite{Browsinglikehuman:luo2025browsing}.  WebExperT outperforms baseline multimodal agents on the Mind2Web benchmark \cite{Browsinglikehuman:luo2025browsing}. WebExperT greatly increases the sophistication of multimodal online agents by combining memory, reflection, and a human-like reasoning style.

\section{Structural DOM Pruning and Local Grounding}
Real-world Website contain ~10,000 to 1,00,000 DOM tokens \cite{zhang2025prune4web}, which is much larger than the LLMs input context windows, and because of this, are unable to use LLM directly to input DOM. \\
\textbf{Prune4Web: DOM Tree Pruning Programming for Web Agent:}\\
Prune4Web \cite{zhang2025prune4web} proposes a scoring-based filter, where the LLM generates a keyword/weight scoring of the generated sub-task plan, and a  Python scorer algorithm uses it over the DOM to generate the top K elements. Thus, the raw DOM is not used as input in LLM directly, and grounder takes only the top k elements, which achieves 25x to 50x token reduction with a reported Element Accuracy of 88.28\% on Mind2Web using GPT-4o as the grounder \cite{zhang2025prune4web}.\\
\textbf{AutoWebGLM: A Large Language Model-based Web Navigating Agent:}\\
AutoWebGLM \cite{lai2024autowebglm} uses the HTML Pruning-based approach for filtering HTML DOM. It limits tree depth and keeps only nodes with actionable descendants. The input given to the model is the task description, simplified HTML, viewport info, and previous action history. Here, the DOM processing is stateless, where the algorithm simplifies the DOM at each stage from scratch \cite{lai2024autowebglm}.\\
\textbf{Agent-E: From Autonomous Web Navigation to Foundational Design Principles in Agentic Systems:}\\
Agent-E \cite{tarsariya2024agente} filters DOM using text-only, input field only, or all fields strategies. It also monitors DOM changes between steps after executing actions, using the browser's MutationObserver API \cite{tarsariya2024agente}. However, the observed change is converted into a natural language description using LLM only and appended to the next prompt, which increases tokens rather than reducing them \cite{tarsariya2024agente}.
None of the above approaches of DOM-diff are wired into the filter stage of an LLM web agent, which leaves an open design space for using DOM Diff observation over the real-world web agents.

\section{DOM State-Change Representations}
This section represents literature that experiments with DOM change observation between two DOM snapshots. \\
\textbf{Beyond Pixels: Exploring DOM Downsampling for LLM-Based Web Agents:}\\
D2Snap \cite{schiepanski2025d2snap} is the first Algorithm designed for downsampling of the DOM rather than truncating or compacting it. It walks the tree in post-order and applies three operations: hierarchical merging of containers, dropping low-relevance text using TextRank, and filtering of attributes below a semantic threshold derived from GPT-4o \cite{schiepanski2025d2snap}. An adaptive wrapper, AdaptiveD2Snap, tunes the three parameters to hit a defined target token budget, forces the DOM size to roughly toward $10^3$ tokens, which matches the token order of grounded GUI screenshots \cite{schiepanski2025d2snap}. The method, used on static DOM snapshots only, does not contain any mechanism for tracking state change across consecutive agent steps \cite{schiepanski2025d2snap}.\\
\textbf{Agentic Compilation: Mitigating the LLM Rerun Crisis for Minimized-Inference-Cost Web Automation:}\\
Agentic Compilation of Web Interfaces \cite{agenticcompilation2025} discusses around DOM difference, prompt compression, and partial state caching in its related work, but it does not implement DOM difference-based system. Its own contribution is a DOM Sanitization Module that achieves around 85\% token reduction through per-step noise removal and attribute cleansing \cite{agenticcompilation2025}.\\
\textbf{RTED: A Robust Algorithm for the Tree Edit Distance:}\\
Tree Edit Distance (TED) \cite{zhangshasha1989ted} algorithm is the canonical algorithm for comparing two tree structures and calculating the minimum-cost sequence of insert/delete/rename operations; RTED \cite{pawlik2011rted} is the current state-of-the-art proposed work with better worst-case performance on skewed trees. These algorithms inform the structural differences between two DOM snapshots, but they were not designed with the real-time token budget of a web agent in mind.\\


\section{Local Grounding Models and their Relevance}
Current Grounding in web agents relies mostly on cloud-based LLMs (Prune4Web) or fine-tuned cross-encoders (Mind2Web), and none of them apply general-purpose sentence embedding methodology over the DOM. \\
\textbf{Sentence-BERT:}
Sentence-BERT \cite{reimers2019sbert} establishes the siamese encoding technique that turns sentences into dense vectors without depending on two sentences only, and shows that late fusion of sparse and dense (embedding) scores usually beats either one alone on retrieval tasks \cite{reimers2019sbert}.\\
\textbf{MiniLM:}
MiniLM \cite{wang2020minilm} separates large Transformer encoders into small task-oriented models. The widely used model \texttt{all-MiniLM-L6-v2} has 6 layers and $\sim$22M parameters, which fits under 100~MB, and runs at roughly 2~ms per sentence on even CPU, which ranks near the top of sentence similarity benchmarks for its size class \cite{wang2020minilm}. These properties matter because a web agent subtask loop has a few seconds to execute grounding with a lower budget, and a heavier encoder can increase inference latency with budget.\\
\textbf{QLoRA and small open grounders:} QLoRA \cite{dettmers2023qlora} showed that on a single consumer GPU with only a small accuracy gap, a 4-bit quantisation of the base model combined with LoRA adapters can fine-tune billion-scale models with full-precision tuning. This allows training a grounder locally rather than relying on a paid hosted model. The Qwen2.5 \cite{qwen25techreport2024} family (and the more recent Qwen3 family) provides permissive, instruction-tuned open-weight bases in the 0.5B to 0.6B range that are realistic targets for such local fine-tuning.



\section{Key Recent Works summary}
Table \ref{Literature:SummaryTable} shows the summary of all recent advancements in the field of GUI automation.

\begin{table}[!h]
\small
\centering
\begin{tabularx}{\textwidth}{X|X|X|c}
\hline
\textbf{Title} & \textbf{Key Contribution} & \textbf{Limitations} & \textbf{Ref.} \\
\hline

Large Language Model-Brained GUI Agents: A Survey &
Gives a clear overview of how GUI agents using LLMs have grown, how they work, and what tools and datasets exist. &
Points out missing areas like stable grounding, better benchmarks, and long-task handling. 
& \cite{zhang2025largelanguagemodelbrainedgui} \\
\hline

A Survey on (M)LLM-Based GUI Agents &
Explains how GUI agents process screens, explore pages, plan actions, and interact. Reviews limits in current methods. &
Notes issues with uneven benchmarks, weak evaluation methods, and poor element detection. 
& \cite{tang2025surveymllmbasedguiagents} \\
\hline

A Survey of WebAgents: Towards Next-Generation AI Agents for Web Automation &
Covers how WebAgents are built, trained, and used. Shows how foundation models make web automation easier. &
Raises concerns about reliability, safety, and the need for stronger training and rules. 
& \cite{ning2025surveywebagentsnextgenerationai} \\
\hline

GUI Agents with Foundation Models: A Comprehensive Survey &
Summarizes FM and LLM-based GUI agent research, giving a simple structure and dataset overview. &
Mentions missing model checkpoints, uneven experiments, and reproducibility problems. 
& \cite{wang2025guiagentsfoundationmodels} \\
\hline

GPT-4V(ision) is a Generalist Web Agent, if Grounded &
Shows GPT-4V can act as a web agent when given screenshots and HTML. Introduces SeeAct with good accuracy. &
Grounding is not perfect; large models and manual setup are still required.
& \cite{SeeAct} \\
\hline

AdaptAgent: Adapting Multimodal Web Agents with Few-Shot Learning &
Uses a few human demos to improve agent performance on new sites; gives noticeable accuracy improvements. &
Depends a lot on which demos are chosen; still struggles with very different websites.
& \cite{FewShotLearning:verma2024adaptagentadaptingmultimodalweb} \\
\hline

LiteWebAgent: The Open-Source Suite for VLM-Based Web-Agent Applications &
Provides an easy, open toolkit with planning, memory, and search modules to build VLM web agents. &
Performance varies based on which model and grounding method are used. 
& \cite{zhang2025litewebagentopensourcesuitevlmbased} \\
\hline

OpenWebAgent: An Open Toolkit to Enable Web Agents on LLMs &
Gives a modular toolkit (HTML parser, action generator, UI) to quickly build web agents; fully open-source. &
Does not provide new accuracy results; depends on how users integrate models.
& \cite{iong2024openwebagent} \\
\hline

Browsing Like Human (WebExperT) &
Adds fast-and-slow reasoning and learning from past mistakes, improving handling of complex tasks. &
More complex design may need extra compute and careful setup.
& \cite{Browsinglikehuman:luo2025browsing} \\
\hline

\end{tabularx}
\caption{Summary of GUI automation recent work}
\label{Literature:SummaryTable}
\end{table}


\section{Critical Research Gaps \& Motivation for Our Work}
Although recent advances, research on multimodal GUI-agents and web automation still GUI automation has important gaps. Recognising these gaps helps motivate work on building a lightweight, DOM-summarising, fallback-enabled agent architecture. Blow are the listed critical gaps:

\begin{enumerate}
    \item \textbf{Unreliable element grounding:} While multimodal like SeeAct \cite{SeeAct} generate action from screenshot images, they rely on mapping of the generated outcome and DOM. SeeAct \cite{SeeAct} provides three grounding strategies, still all of them struggle with grounding and drop agent performance due to element grounding limitation.SeeAct \cite{SeeAct} motivates us to work on element grounding, where instead of relying only on the key pair comparison, we directly ask the LLM for the exact grounding value.

    \item \textbf{Privacy concern and Human-in-Loop:} GUI agents become more autonomous and powerful, which raises concerns about security, privacy, and safe execution. Survey \cite{tang2025surveymllmbasedguiagents} \cite{ning2025surveywebagentsnextgenerationai} call out deficiencies in safety-aware execution control and benchmark standardisation.

    \item \textbf{Large and complex DOMs, missing js, css:} Real-world websites contain thousands of DOM elements, where model has their input limitation beyond which they won't work \cite{wang2025guiagentsfoundationmodels}. Model does not look at JS or CSS while inference, thus loses the context where the user prompt contains context about visual values \cite{wang2025guiagentsfoundationmodels}. 
    
    \item \textbf{Cost and Token Inefficiency with Large Models:} Using large multimodal models (like GPT-4V) along with full HTML/DOM inputs and images is resource-intensive and often infeasible for small-scale or low-resource environments, which raises concern regarding cost and input limitations \cite{tang2025surveymllmbasedguiagents}.

    \item \textbf{Stateless DOM processing:} All recent DOM reduction based WebAutomation approaches \cite{zhang2025prune4web} \cite{lai2024autowebglm} \cite{schiepanski2025d2snap}  process each subtask DOM independently. In most multi-step tasks, the page only changes in a way where interaction or knowledge is required. Agent-E's change observer \cite{tarsariya2024agente} is the only system that looks at DOM change across steps, but it uses that DOM Difference to \textit{describe} what happened (action feedback) rather than to \textit{reduce} what is sent to the grounder, so it increases tokens instead of saving them.

    \item \textbf{Privacy and safety bolted on post-hoc:} Surveys \cite{tang2025surveymllmbasedguiagents} \cite{ning2025surveywebagentsnextgenerationai} state on safety-aware execution and privacy as open concerns, but current approaches, which address DOM reduction and grounding, that do not design or focued on  privacy at the filter/grounder interface level. Personally Identifiable Information (PII) in form fields, password placeholders, and identification numbers end up in the LLM context by default, and human-in-the-loop gates are added as a final confirm step rather than as a risk-aware layer that runs before grounding.
    
\end{enumerate}
